% Integrated 4 September 2026; exact in-project provenance is recorded below.
% The algebra, order, lattice, shell, centralizer, and cutoff are retained exactly.
\section{Exact quaternion bridge: order, shell, and hyperbolic conjugacy classes}\label{koide-quat-exact-quaternion-bridge-order-shell-and-hyperbolic-conjugacy-classes}

The proposed quaternionic realization is correct. The displayed lattice is an order in an indefinite rational quaternion division algebra of discriminant \(6\). The parity-restricted ternary shell is exactly the reduced-trace-\(4\), reduced-norm-\(1\) locus in that order. Every resulting element is primitive hyperbolic, with eigenvalues \(2\pm\sqrt3\) and translation length \(2\log(2+\sqrt3)\).

The centralizer statement is correct when written intrinsically as

\[
 \mathcal O\cap\mathbb Q(\gamma)
 =\mathbb Z[\gamma-2]\cong\mathbb Z[\sqrt3].
\]

The entire integral shell is a finite union of conjugacy classes under \(\Gamma=\mathcal O^1\). This gives an exact reduction of the Koide count to a finite sum of sharp, non-radial sector counts inside primitive hyperbolic conjugacy classes of a cocompact arithmetic Fuchsian group.

It does \textbf{not} prove the open estimate

\[
 R(X)-C_0X=O_\varepsilon(X^{1/2+\varepsilon}).
\]

Rather, it converts the count exactly into a sharp sector problem on a finite union of hyperbolic conjugacy classes. This places it beside, but does not prove it equivalent to, the standard hyperbolic conjugacy-class lattice-point problem. Even for the radial Huber cutoff, the general pointwise theorem gives an \(O(X^{2/3})\) remainder after all exceptional spectral terms are included, while the \(O(X^{1/2+\varepsilon})\) remainder is conjectural. The Koide cutoff is a different sharp polyhedral-sector height, so the standard radial theorem cannot be imported without a new transfer or trace-formula argument.

\subsection{Quaternion algebra and reduced invariants}\label{koide-quat-quaternion-algebra-and-reduced-invariants}

Let

\[
 B=(18,-6)_{\mathbb Q}
 =\mathbb Q\langle I,J\rangle/
 (I^2-18,\ J^2+6,\ IJ+JI),
 \qquad K=IJ.
\]

Then

\[
 K^2=-I^2J^2=108.
\]

For

\[
 x=x_0+x_1I+x_2J+x_3K,
\]

the standard involution, reduced trace, and reduced norm are

\[
 \bar x=x_0-x_1I-x_2J-x_3K,
 \qquad \operatorname{trd}(x)=x+\bar x=2x_0,
\]

and

\begin{equation}\label{koide-quat-eq-reduced-norm}
 \operatorname{nrd}(x)=x\bar x
 =x_0^2-18x_1^2+6x_2^2-108x_3^2.
\end{equation}

All signs in \eqref{koide-quat-eq-reduced-norm} are forced by \(\operatorname{nrd}(I)=-18\), \(\operatorname{nrd}(J)=6\), and \(\operatorname{nrd}(K)=-108\).

\subsection{The proposed lattice is an order}\label{koide-quat-the-proposed-lattice-is-an-order}

Put

\[
 E=\frac{I+J}{2},\qquad
 F_0=\frac{I-J}{2},\qquad
 G=\frac K2,
\]

and let

\begin{equation}\label{koide-quat-eq-order-lattice}
 \mathcal O=\mathbb Z1\oplus\mathbb ZE
             \oplus\mathbb ZF_0\oplus\mathbb ZG.
\end{equation}

The displayed elements are a rational basis of \(B\), because

\[
 I=E+F_0,\qquad J=E-F_0,\qquad K=2G.
\]

Direct multiplication gives the complete nontrivial table

\begin{equation}\label{koide-quat-eq-multiplication-table}
\begin{array}{c|ccc}
   &E&F_0&G\\ \hline
 E&3&6-G&6E-3F_0\\
 F_0&6+G&3&3E-6F_0\\
 G&-6E+3F_0&-3E+6F_0&27.
\end{array}
\end{equation}

Every structure constant in \eqref{koide-quat-eq-multiplication-table} is integral. Thus \eqref{koide-quat-eq-order-lattice} is a unital subring and a full rank-four \(\mathbb Z\)-lattice in \(B\), hence an order. In particular, the unital ring sometimes denoted \(\mathbb Z[1,E,F_0,G]\) is exactly the direct sum in \eqref{koide-quat-eq-order-lattice}.

An arbitrary order element is

\begin{equation}\label{koide-quat-eq-order-coordinates}
 x=r+sE+tF_0+uG
  =r+\frac{s+t}{2}I+\frac{s-t}{2}J+\frac u2K
 \qquad(r,s,t,u\in\mathbb Z),
\end{equation}

and \eqref{koide-quat-eq-reduced-norm} gives

\begin{equation}\label{koide-quat-eq-order-trace-norm}
 \operatorname{trd}(x)=2r,
 \qquad
 \operatorname{nrd}(x)
 =r^2-3s^2-12st-3t^2-27u^2.
\end{equation}

The reduced-trace pairing on the basis \((1,E,F_0,G)\) is

\[
 \bigl(\operatorname{trd}(b_ib_j)\bigr)_{i,j}
 =
 \begin{pmatrix}
  2&0&0&0\\
  0&6&12&0\\
  0&12&6&0\\
  0&0&0&54
 \end{pmatrix},
\]

with determinant

\begin{equation}\label{koide-quat-eq-order-discriminant-determinant}
 -11664=-108^2.
\end{equation}

Thus the reduced discriminant of this order is \(108\), not \(6\). The algebra has discriminant \(6\), proved below, so \(\mathcal O\) is nonmaximal and has index \(18\) in a maximal order containing it. This index is not an Eichler level: the order is nonmaximal at the ramified primes \(2\) and \(3\), so it is not an Eichler order in the standard sense. Any optimal-embedding or trace-formula result that assumes a maximal order must account for this exact nonmaximal order instead of replacing \(\mathcal O\).

\subsection{Trace-zero lattice and exact shell bijection}\label{koide-quat-trace-zero-lattice-and-exact-shell-bijection}

Equation~\eqref{koide-quat-eq-order-trace-norm} gives

\[
 \mathcal O^0
 :=\{x\in\mathcal O:\operatorname{trd}(x)=0\}
 =\mathbb ZE\oplus\mathbb ZF_0\oplus\mathbb ZG.
\]

Writing

\[
 q=s+t,\qquad a=s-t,\qquad n=u
\]

gives the exact description

\begin{equation}\label{koide-quat-eq-trace-zero-lattice}
 \boxed{
 \mathcal O^0=
 \left\{
 \frac{qI+aJ+nK}{2}:
 a,q,n\in\mathbb Z,\ a\equiv q\pmod2
 \right\}.}
\end{equation}

The inverse formulas are

\[
 s=\frac{q+a}{2},\qquad
 t=\frac{q-a}{2},\qquad
 u=n,
\]

which are integral exactly when \(a\equiv q\pmod2\).

Define

\[
 F(a,q,n)=a^2-3q^2-18n^2,
 \qquad
 v(a,q,n)=\frac{qI+aJ+nK}{2},
\]

and

\begin{equation}\label{koide-quat-eq-gamma-parametrization}
 \gamma(a,q,n)=2+v(a,q,n).
\end{equation}

Distinct pure quaternion basis elements anticommute, so

\begin{equation}\label{koide-quat-eq-pure-square}
 v(a,q,n)^2
 =\frac{18q^2-6a^2+108n^2}{4}
 =-\frac32F(a,q,n),
\end{equation}

and

\[
 \operatorname{nrd}(v)=\frac32F(a,q,n).
\]

Since \(v\) has reduced trace zero,

\begin{equation}\label{koide-quat-eq-gamma-trace-norm}
 \operatorname{trd}(\gamma)=4,
 \qquad
 \operatorname{nrd}(\gamma)
 =4+\operatorname{nrd}(v)
 =4+\frac32F(a,q,n).
\end{equation}

Therefore \eqref{koide-quat-eq-gamma-parametrization} defines a bijection

\begin{equation}\label{koide-quat-eq-shell-bijection}
 \boxed{
 \left\{(a,q,n)\in\mathbb Z^3:
 F(a,q,n)=-2,\ a\equiv q\pmod2\right\}
 \xrightarrow{\ \sim\ }
 \left\{\gamma\in\mathcal O:
 \operatorname{trd}(\gamma)=4,\
 \operatorname{nrd}(\gamma)=1\right\}.}
\end{equation}

The inverse sends

\[
 \gamma=2+sE+tF_0+uG
\]

to

\[
 (a,q,n)=(s-t,s+t,u).
\]

There is no hidden congruence loss. On the integral equation \(F=-2\), reduction modulo \(2\) gives

\[
 a^2-q^2\equiv0\pmod2,
\]

so \(a\equiv q\pmod2\) is automatic. The parity condition in \eqref{koide-quat-eq-shell-bijection} records the order lattice exactly but removes no integral shell point.

Finally,

\begin{equation}\label{koide-quat-eq-gamma-quadratic}
 (\gamma-2)^2=3,\qquad
 \gamma^2-4\gamma+1=0,\qquad
 \gamma^{-1}=4-\gamma=2-v\in\mathcal O.
\end{equation}

Thus every element on the right of \eqref{koide-quat-eq-shell-bijection} lies in \(\Gamma=\mathcal O^1\).

\subsection{\texorpdfstring{Division over \(\mathbb Q\) and splitting over \(\mathbb R\)}{Division over \textbackslash mathbb Q and splitting over \textbackslash mathbb R}}\label{koide-quat-division-over-mathbb-q-and-splitting-over-mathbb-r}

Square factors may be removed from a Hilbert-symbol presentation, so

\[
 (18,-6)_{\mathbb Q}\simeq(2,-6)_{\mathbb Q}.
\]

At \(p=3\), write

\[
 18=3^2\cdot2,\qquad -6=3\cdot(-2).
\]

The odd-prime Hilbert-symbol formula gives

\begin{equation}\label{koide-quat-eq-hilbert-three}
 (18,-6)_3
 =(-1)^{2\cdot1\cdot(3-1)/2}
   \left(\frac2{3}\right)
   \left(\frac{-2}{3}\right)^2
 =-1.
\end{equation}

Thus \(B\otimes\mathbb Q_3\) is division, and \(B\) is a division algebra over \(\mathbb Q\).

At \(p=2\), write

\[
 18=2\cdot9,\qquad -6=2\cdot(-3).
\]

The \(2\)-adic formula

\[
 (2^\alpha u,2^\beta w)_2
 =(-1)^{
 ((u-1)/2)((w-1)/2)
 +\alpha(w^2-1)/8+\beta(u^2-1)/8}
\]

gives

\begin{equation}\label{koide-quat-eq-hilbert-two}
 (18,-6)_2=-1.
\end{equation}

At every odd prime other than \(3\), both valuations are zero and the Hilbert symbol is \(+1\). At the real place, \(18>0\), so the real Hilbert symbol is \(+1\). Consequently

\begin{equation}\label{koide-quat-eq-ramification-discriminant}
 \operatorname{Ram}(B)=\{2,3\},\qquad
 \operatorname{disc}(B)=6,
\end{equation}

and

\[
 B\otimes_{\mathbb Q}\mathbb R\simeq M_2(\mathbb R).
\]

An explicit real splitting is

\begin{equation}\label{koide-quat-eq-real-splitting}
 \rho(I)=
 \begin{pmatrix}3\sqrt2&0\\0&-3\sqrt2\end{pmatrix},
 \qquad
 \rho(J)=
 \begin{pmatrix}0&6\\-1&0\end{pmatrix}.
\end{equation}

These matrices square to \(18\) and \(-6\) and anticommute. Therefore

\[
 \rho(K)=
 \begin{pmatrix}0&18\sqrt2\\3\sqrt2&0\end{pmatrix}.
\]

For the element associated with \((a,q,n)\),

\begin{equation}\label{koide-quat-eq-gamma-matrix}
 \rho(\gamma)=
 \begin{pmatrix}
 2+\frac{3\sqrt2}{2}q&3a+9\sqrt2\,n\\
 -\frac12a+\frac{3\sqrt2}{2}n&
 2-\frac{3\sqrt2}{2}q
 \end{pmatrix}.
\end{equation}

Its determinant is the reduced norm in \eqref{koide-quat-eq-gamma-trace-norm}.

Let \(\mathcal M\) be a maximal overorder containing \(\mathcal O\).
Since \([\mathcal M:\mathcal O]=18\), the finite abelian group
\(\mathcal M/\mathcal O\) has order \(18\).  Lagrange's theorem therefore
shows that multiplication by \(18\) annihilates every element of this
additive quotient; equivalently, the ideal \((18)=18\mathbb Z\) annihilates
\(\mathcal M/\mathcal O\), and \(18\mathcal M\subset\mathcal O\).  Hence the
principal congruence subgroup
\[
 \mathcal M^1(18\mathcal M)
 :=\{u\in\mathcal M^1:u-1\in18\mathcal M\}
\]
satisfies the exact inclusions
\[
 \mathcal M^1(18\mathcal M)\subset\mathcal O^1\subset\mathcal M^1.
\]
Thus \(\mathcal O^1\) has finite index in the norm-one group of a maximal
order and contains a principal congruence subgroup.  This is the required
finite-index and congruence bridge; it does not replace the local type of
\(\mathcal O\) by the local type of \(\mathcal M\).

The image

\[
 \bar\Gamma=\rho(\mathcal O^1)/\{\pm1\}
 \subset PSL_2(\mathbb R)
\]

is a discrete congruence arithmetic Fuchsian group. Since the defining
rational quaternion algebra is division,
\(\bar\Gamma\backslash\mathbb H\) is compact; see
\cite[Lemma~38.1.2, Proposition~38.1.3, and Definition~38.7.2]{Voight_2021}.

\subsection{Centralizer, optimal quadratic order, and primitivity}\label{koide-quat-centralizer-optimal-quadratic-order-and-primitivity}

Fix a shell point and set

\[
 v=\gamma-2.
\]

Equation~\eqref{koide-quat-eq-gamma-quadratic} gives \(v^2=3\), hence

\begin{equation}\label{koide-quat-eq-quadratic-field}
 \mathbb Q(\gamma)=\mathbb Q(v)\cong\mathbb Q(\sqrt3).
\end{equation}

Since \(\gamma\) is noncentral, the centralizer theorem for quaternion algebras gives

\begin{equation}\label{koide-quat-eq-algebra-centralizer}
 C_B(\gamma)=\mathbb Q(\gamma).
\end{equation}

Moreover,

\[
 \mathbb Z[v]\subseteq\mathcal O\cap\mathbb Q(v).
\]

Every element of \(\mathcal O\) is integral over \(\mathbb Z\), so every element of the intersection is an algebraic integer in \(\mathbb Q(\sqrt3)\). Because

\[
 \mathcal O_{\mathbb Q(\sqrt3)}=\mathbb Z[\sqrt3]
\]

and \(\mathbb Z[v]\) is already isomorphic to this maximal quadratic order, both inclusions force

\begin{equation}\label{koide-quat-eq-optimal-intersection}
 \boxed{
 \mathcal O\cap\mathbb Q(\gamma)
 =\mathbb Z[v]=\mathbb Z[\gamma-2]
 \cong\mathbb Z[\sqrt3].}
\end{equation}

Thus each shell point gives an optimal embedding

\begin{equation}\label{koide-quat-eq-optimal-embedding}
 \iota_\gamma:\mathbb Z[\sqrt3]\hookrightarrow\mathcal O,
 \qquad \sqrt3\longmapsto\gamma-2.
\end{equation}

Conversely, every embedding \eqref{koide-quat-eq-optimal-embedding} gives a trace-\(4\), norm-\(1\) element \(2+\iota(\sqrt3)\), hence one integral ternary point. The shell is exactly the set of these optimal embeddings before quotienting by conjugation.

The norm-one units of \(\mathbb Z[\sqrt3]\) are

\begin{equation}\label{koide-quat-eq-norm-one-units}
 \{\pm(2+\sqrt3)^k:k\in\mathbb Z\}.
\end{equation}

For completeness, \(2+\sqrt3\) is the smallest positive norm-one unit greater than \(1\). If \(x+y\sqrt3>1\) has norm \(1\) and \(y>0\), then \(y\ge1\), \(x^2=1+3y^2\), and \(x+y\sqrt3\ge2+\sqrt3\). Multiplication by a suitable power of \(2-\sqrt3\) reduces any positive norm-one unit to \([1,2+\sqrt3)\), proving \eqref{koide-quat-eq-norm-one-units}.

Equations~\eqref{koide-quat-eq-algebra-centralizer}--\eqref{koide-quat-eq-norm-one-units} imply

\begin{equation}\label{koide-quat-eq-group-centralizer}
 C_\Gamma(\gamma)=\{\pm\gamma^k:k\in\mathbb Z\}.
\end{equation}

Hence \(\gamma\) is primitive in \(\Gamma\) and its image is primitive in \(\bar\Gamma\). Indeed, a root of \(\gamma\), or a root of its class modulo the center, would commute with \(\gamma\) and therefore be a unit from \eqref{koide-quat-eq-norm-one-units}; the exponent equation \(rk=1\) has no solution with \(k\ge2\).

The reduced characteristic polynomial is

\[
 T^2-4T+1,
\]

so the real eigenvalues are

\begin{equation}\label{koide-quat-eq-eigenvalues}
 2+\sqrt3,\qquad 2-\sqrt3.
\end{equation}

For a hyperbolic element of \(PSL_2(\mathbb R)\),

\[
 2\cosh(\ell/2)=|\operatorname{tr}(\rho(\gamma))|=4.
\]

Thus every shell element has primitive translation length

\begin{equation}\label{koide-quat-eq-translation-length}
 \boxed{
 \ell_0=2\operatorname{arcosh}(2)
       =2\log(2+\sqrt3).}
\end{equation}

\subsection{Integral adjoint action and finitely many classes}\label{koide-quat-integral-adjoint-action-and-finitely-many-classes}

Let

\[
 \mathscr S
 =\{\gamma\in\Gamma:\operatorname{trd}(\gamma)=4\}.
\]

By \eqref{koide-quat-eq-shell-bijection}, \(\mathscr S\) is exactly the entire integral ternary shell. By \eqref{koide-quat-eq-translation-length}, every element of \(\mathscr S\) is hyperbolic with translation length \(\ell_0\).

Put

\[
 L=\{(a,q,n)\in\mathbb Z^3:a\equiv q\pmod2\}.
\]

Conjugation preserves \(\mathcal O\), \(\mathcal O^0\), and reduced norm. Because

\[
 \operatorname{nrd}(v(a,q,n))=\frac32F(a,q,n),
\]

it gives an exact orthogonal representation on the actual parity lattice

\begin{equation}\label{koide-quat-eq-adjoint-representation}
 \operatorname{Ad}:\Gamma/\{\pm1\}
 \hookrightarrow SO^+(L,F).
\end{equation}

The kernel before dividing by the center is \(\Gamma\cap\mathbb Q^\times=\{\pm1\}\). No action on the larger lattice \(\mathbb Z^3\), and no finite-index equality with its full integral orthogonal group, is asserted.

There are only finitely many \(\Gamma\)-conjugacy classes in \(\mathscr S\). Choose a compact fundamental set \(D\subset\mathbb H\) for \(\bar\Gamma\), and fix \(z_0\in\mathbb H\). For each conjugacy class, conjugate a representative so that its invariant axis meets \(D\). If \(z\in D\) lies on that axis, then

\[
 d(z,\gamma z)=\ell_0.
\]

Writing

\[
 R_D=\max_{w\in D}d(z_0,w),
\]

the triangle inequality gives

\[
 d(z_0,\gamma z_0)\le2R_D+\ell_0.
\]

Proper discontinuity leaves only finitely many elements of \(\bar\Gamma\) moving \(z_0\) by at most this amount. Hence only finitely many \(\bar\Gamma\)-classes occur. This also proves finiteness of the corresponding \(\Gamma\)-classes: if two images are conjugate in \(\bar\Gamma\), lifting a conjugator gives

\[
 \delta\gamma_1\delta^{-1}=\pm\gamma_2.
\]

Both sides before the possible sign have reduced trace \(4\); the minus sign would have reduced trace \(-4\), so it is impossible.

Choose representatives \(\gamma_1,\ldots,\gamma_h\). Then

\begin{equation}\label{koide-quat-eq-shell-class-decomposition}
 \boxed{
 \mathscr S=
 \coprod_{j=1}^h
 \{\delta\gamma_j\delta^{-1}:
 [\delta]\in\Gamma/C_\Gamma(\gamma_j)\},
 \qquad
 C_\Gamma(\gamma_j)=\{\pm\gamma_j^k:k\in\mathbb Z\}.}
\end{equation}

This proves finiteness but does not compute \(h\) or a representative list. Those finite arithmetic data are needed by a class-by-class spectral formula whose coefficients contain geodesic periods.

\subsection{Exact reformulation of the Koide chamber count}\label{koide-quat-exact-reformulation-of-the-koide-chamber-count}

For \(\gamma\in\mathscr S\), let

\[
 c(\gamma)=(a(\gamma),q(\gamma),n(\gamma))
\]

be the inverse coordinate map in
\cref{koide-quat-trace-zero-lattice-and-exact-shell-bijection}. Let

\[
 \mathcal C=\{(a,q,n)\in\mathbb R^3:
 a>0,\ q\ge0,\ n\ge0,\ a+4n>q\},
\]

and retain the canonical height

\[
 H(a,q,n)=4a+17n.
\]

The integral bijections in
\cref{koide-count-thm-coordinate-bijection,koide-param-thm:integral-parametrization}
give

\begin{equation}\label{koide-quat-eq-koide-class-count}
 \boxed{
 R(X)=\sum_{j=1}^h
 \#\left\{
 [\delta]\in\Gamma/C_\Gamma(\gamma_j):
 \begin{array}{l}
 c(\delta\gamma_j\delta^{-1})\in\mathcal C,\\
 H(c(\delta\gamma_j\delta^{-1}))\le X
 \end{array}
 \right\}.}
\end{equation}

There is no hidden finite-to-one factor in
\eqref{koide-quat-eq-koide-class-count}: cosets by the centralizer
parametrize the elements of each class exactly once, and
\eqref{koide-quat-eq-shell-bijection} parametrizes shell elements by triples
exactly once.  Put \(C_j=C_\Gamma(\gamma_j)\).  Here
\(\Gamma/C_j\) is the set of left cosets \(\delta C_j\), not the set of
right cosets \(C_j\delta\), and the exact orbit bijection is

\[
 \Phi_j:\Gamma/C_j\longrightarrow
 \{\delta\gamma_j\delta^{-1}:\delta\in\Gamma\},
 \qquad
 \delta C_j\longmapsto\delta\gamma_j\delta^{-1}.
\]

Indeed, right multiplication of \(\delta\) by an element of \(C_j\) does
not change its conjugate, so \(\Phi_j\) is well defined and surjective; and

\[
 \Phi_j(\delta_1C_j)=\Phi_j(\delta_2C_j)
 \iff \delta_2^{-1}\delta_1\in C_j
 \iff \delta_1C_j=\delta_2C_j,
\]

so it is injective.  If \(\pi:\Gamma\to\bar\Gamma\) is the quotient map,
then

\[
 \Gamma/C_j
 \cong
 \bar\Gamma/C_{\bar\Gamma}(\bar\gamma_j)
\]

via \(\delta C_j\mapsto\pi(\delta)C_{\bar\Gamma}(\bar\gamma_j)\).
To prove the required equality of stabilizers, observe first that
\(\pi(C_j)\subseteq C_{\bar\Gamma}(\bar\gamma_j)\).  Conversely, if
\(\pi(\delta)\) centralizes \(\bar\gamma_j\), then
\(\delta\gamma_j\delta^{-1}=\pm\gamma_j\).  Reduced-trace invariance and
\(\operatorname{trd}(\gamma_j)=4\) exclude the minus sign, so
\(\delta\in C_j\).  Thus
\(\pi(C_j)=C_{\bar\Gamma}(\bar\gamma_j)\); together with
\(\ker\pi=\{\pm1\}\subset C_j\), this proves the displayed bijection and
shows that passage to \(PSL_2\) introduces no factor \(2\). Positive trace
selects the lift sign, but the chamber is not invariant under
\(\gamma\mapsto\gamma^{-1}\); inverse classes and their orientations must
still be retained separately.

Under the splitting \eqref{koide-quat-eq-real-splitting},

\begin{equation}\label{koide-quat-eq-frobenius-norm}
 \|\rho(\gamma)\|_{\mathrm F}^2
 =8+9q^2+\frac{37}{4}a^2
   +\frac{105\sqrt2}{2}an+\frac{333}{2}n^2.
\end{equation}

For \(i\in\mathbb H\),

\[
 2\cosh d(i,\rho(\gamma)i)=\|\rho(\gamma)\|_{\mathrm F}^2.
\]

Since \(\sinh(\ell_0/2)=\sqrt3\), the standard Huber parameter is

\begin{equation}\label{koide-quat-eq-huber-parameter}
 X_i(\gamma)^2
 =\frac{\|\rho(\gamma)\|_{\mathrm F}^2-2}{12}.
\end{equation}

The shell equation simplifies \eqref{koide-quat-eq-frobenius-norm} exactly:

\[
 \|\rho(\gamma)\|_{\mathrm F}^2
 =14+\frac{(7a+15\sqrt2\,n)^2}{4},
\]

and hence

\begin{equation}\label{koide-quat-eq-radial-height}
 \boxed{
 X_i(\gamma)^2
 =1+\frac{(7a+15\sqrt2\,n)^2}{48}.}
\end{equation}

Thus, in the positive chamber, the radial cutoff has the exact irrational linear height \(7a+15\sqrt2\,n\). The Koide height \(H=4a+17n\) is a different linear height. The two sizes are comparable:

\[
 a\le\frac H4,\qquad
 n\le\frac H{17},\qquad
 0\le q<a+4n\le\frac{33}{68}H,
\]

and conversely

\[
 H\le\sqrt{305}\sqrt{a^2+q^2+n^2}.
\]

Norm equivalence in the trace-zero space then compares \(H\) coarsely with \(X_i\). Such comparability identifies the growth scale but does not transfer a main coefficient or a remainder between the sharply bounded regions.

\subsection{Consequence for the open square-root discrepancy}\label{koide-quat-consequence-for-the-open-square-root-discrepancy}

The quaternion construction proves:

\begin{enumerate}
\def\labelenumi{\arabic{enumi}.}
\item
  the ternary shell is exactly a finite union of primitive hyperbolic conjugacy classes in one explicit cocompact arithmetic Fuchsian group;
\item
  \(\log(2+\sqrt3)=\ell_0/2\), so its occurrence in the main constant is tied to the primitive translation length;
\item
  the conjectural exponent \(1/2\) is the same exponent proposed for the standard hyperbolic lattice-point problem in conjugacy classes, giving a precise analytic analogy but not an equivalence of the two cutoff problems.
\end{enumerate}

The third item classifies the remaining analytic problem; it does not prove the
error estimate. Chatzakos--Petridis~\cite[Theorem~1.1]{ChatzakosPetridis_2016}
give

\[
 E(\mathcal H,X;z)=O(X^{2/3})
\]

for the standard radial count after the main expression, including every small-eigenvalue term, is subtracted. Their Conjecture 5.7 proposes

\[
 E(\mathcal H,X;z)=O_\varepsilon(X^{1/2+\varepsilon}).
\]

The cutoff in \eqref{koide-quat-eq-koide-class-count} uses both a sharp chamber and \(4a+17n\), not the radial height \(7a+15\sqrt2\,n\) in \eqref{koide-quat-eq-radial-height}. Their symmetric difference can have main-order size, so even the conjectured radial bound would not by itself prove the Koide bound. Therefore the status of

\[
 R(X)-C_0X=O_\varepsilon(X^{1/2+\varepsilon})
\]

remains \textbf{OPEN}. Closing it requires a spectral or automorphic argument for the exact finite sum \eqref{koide-quat-eq-koide-class-count}, with its exact sector and height.

\subsection{Missing mathematics before an error theorem can be claimed}\label{koide-quat-missing-mathematics-before-an-error-theorem-can-be-claimed}

The bridge does not supply the following results.

\begin{enumerate}
\def\labelenumi{\arabic{enumi}.}
\item
  \textbf{Class data.} Determine \(h\), representatives \(\gamma_j\), and orientation/inversion conventions. Finiteness does not compute the periods in a spectral expansion.
\item
  \textbf{Exact nonmaximal order.} The order has reduced discriminant \(108\), is nonmaximal at the ramified primes, and is not an Eichler order in the standard sense. A result for maximal or specified Eichler orders must be checked at this order or transferred through finite-index groups with all coset multiplicities retained. Index and discriminant alone do not determine the precise local order type required by a trace formula.
\item
  \textbf{Central quotient.} Track the passage from \(\Gamma\subset SL_2(\mathbb R)\) to \(\bar\Gamma\subset PSL_2(\mathbb R)\), including \(\pm1\), inverse classes, and possible orbifold torsion.
\item
  \textbf{Exact cutoff.} Prove a trace formula or effective orbital theorem for all four inequalities defining \(\mathcal C\) and for \(4a+17n\le X\). The radial Huber cutoff is different.
\item
  \textbf{Sharp-boundary control.} Remove smoothing with an explicit estimate at the required scale. The two Pell faces contribute only \(O(\log X)\), as already proved canonically, but this alone does not control every smoothing strip adjacent to all chamber walls.
\item
  \textbf{Exceptional spectrum.} Once an appropriate sector spectral expansion is constructed, one must determine whether it contains \(X^{s_j}\) terms for \(\lambda_j=s_j(1-s_j)<1/4\), weighted by closed-geodesic periods. To subtract only \(C_0X\) with a \(O(X^{1/2+\varepsilon})\) remainder, every relevant exceptional term must be shown to vanish or cancel, or it must be displayed and subtracted. No such vanishing is proved here.
\item
  \textbf{Pointwise square-root estimate.} After secondary terms and the exact sector are handled, the needed pointwise \(1/2+\varepsilon\) bound remains at the conjectural strength of the classical problem. Mean-square results do not give it for every \(X\).
\item
  \textbf{Main coefficient.} A spectral derivation must reproduce \[
  C_0=\frac{9}{2\pi^2}\log(2+\sqrt3)
      \log\!\left(\frac3{2\sqrt2}\right)
  \] using the same lattice, chamber, strict walls, height, and Haar measure. The existing local-density proof fixes this constant but does not calibrate a new trace formula automatically.
\item
  \textbf{Transfer between the two exact heights.} Formula~\eqref{koide-quat-eq-radial-height} already gives the exact standard radial parameter. What is missing is a remainder-level transfer between its height \(7a+15\sqrt2\,n\) and the canonical height \(4a+17n\), together with the chamber walls. Coarse norm comparability does not preserve a discrepancy estimate.
\end{enumerate}

These are concrete mathematical tasks. None invalidates the quaternion bridge; together they explain why it does not close the open claim.

\subsection{Reproducibility and references}\label{koide-quat-reproducibility-and-references}

The canonical quaternion replay is
\path{supporting_materials/workbench/evidence/koide_exact_checkpoint_20260904/quaternion/verify_quaternion_bridge.py}
(SHA-256
\path{6675ba2fc4fa8c2b10e14558e412178b6b5b579c9c96afca78e1db9bcdf5c6b4}).
Its recorded successful output is
\path{supporting_materials/workbench/evidence/koide_exact_checkpoint_20260904/quaternion/verification.log}
(SHA-256
\path{66853c6191a000c7c827aacd1e0e33dcf34dbe3dd61d746f66fa4d805f34c91e}),
and its local manifest is
\path{supporting_materials/workbench/evidence/koide_exact_checkpoint_20260904/quaternion/SHA256SUMS.txt}
(SHA-256
\path{f9d044db0d4b003e7aef7649b3394e6b299c09a6b2eb7adca095f8f9100ba623}).
The checkpoint-wide manifest is
\path{supporting_materials/workbench/evidence/koide_exact_checkpoint_20260904/SHA256SUMS}
(SHA-256
\path{8d108eb928e5d07093badde169f443a4a38b6fcf3eb59d5192bdd33ecef8d7a4}).
The replay checks the basis multiplication table; the trace, norm, and
coordinate formulas for
\((r,s,t,u)\in\{-4,-3,\ldots,4\}^4\); the shell identities for
\((a,q,n)\in\{-24,-23,\ldots,24\}^2\times\{-8,-7,\ldots,8\}\), containing
\(48\) shell points; the
inverse identity on those points; the radial-height identity on that box;
and the Hilbert-symbol formulas at \(2\) and \(3\).  These finite ranges are
consistency checks, not proofs outside the stated boxes.  The exact canonical
sources used for the counting statements imported into this bridge are:

\begin{itemize}
\item
  \path{supporting_materials/workbench/research/koide_counting_repair.tex},
  SHA-256
  \path{965b0a3dca8bfe4f22b9e9208368fa04697070fd25bd06bd67f2481d17cb9d9d};
\item
  \path{supporting_materials/workbench/research/koide_exact_parametrization.tex},
  SHA-256
  \path{a457709dcce0fce0774eb9326b84b00d5d8144c2087f480da704974e499bcdd3};
\item
  \path{supporting_materials/workbench/research/koide_exact_census.tex},
  SHA-256
  \path{18ed5968b7e67d9a0b23d8675b40c37ba46170ad29efdbb9ce10ed632c99600a}.
\end{itemize}

The bounded coordinate-bijection replay, whose scope does not include the
quaternion multiplication table, Hilbert symbols, centralizer, conjugacy
finiteness, or spectral conclusions, is
\path{supporting_materials/workbench/reproduction/t1_koide_exact/verify_koide_narain_root_density.py}
(SHA-256
\path{e4bf2982b26b07add6057b5d2d5c771a52ef52c629b5ad5f0aab24e440e41986}).
Its stored height-\(200000\) output is
\path{supporting_materials/workbench/reproduction/t1_koide_exact/koide_verification_200000.json}
(SHA-256
\path{b907f0ee5bdd3ad80622b669722d218b79a11a5baab0d4571e141936100706b7});
the exact invocation and scope are stated in
\path{supporting_materials/workbench/reproduction/t1_koide_exact/README.md}
(SHA-256
\path{5e9cdcb9e67954c37d8fde69e13549172aa03245674b358e4a1d25f8d9d3bff0}).
The quaternion identities and deductions in this section are proved in the
displayed algebraic calculations themselves, not inferred from that bounded
replay.

References:

\begin{itemize}
\item
  John Voight, \emph{Quaternion Algebras}, Graduate Texts in Mathematics 288,
  especially Section 10.1, Lemma 38.1.2, and Proposition 38.1.3;
  \href{https://doi.org/10.1007/978-3-030-56694-4}{DOI}.
\item
  Dimitrios Chatzakos and Yiannis N. Petridis,
  \emph{The hyperbolic lattice point problem in conjugacy classes},
  arXiv:1504.01307v2, Theorem 1.1 and Conjecture 5.7;
  \href{https://arxiv.org/abs/1504.01307}{arXiv}.
\item
  The exact Koide chamber, height, main term, and Pell-face estimates are in
  Sections~\ref{koide-count-sec-main}--\ref{koide-count-sec-census-1e8}
  of this workbench.
\end{itemize}
